 % !TEX program = lualatex   
\documentclass[12pt, a4paper]{book}

\raggedbottom

% Packages
\usepackage{amsmath, amssymb, amsthm}
\usepackage{geometry}
\usepackage{hyperref}
\usepackage{enumitem}
\usepackage{unicode-math}
\usepackage{emoji}
\usepackage{xcolor}
\usepackage{tikz-cd}
\usepackage{mathtools}
\usepackage{pifont}
\usepackage{gensymb}

\geometry{margin=1in}

% Theorem environments
\theoremstyle{definition}
\newtheorem{definition}{Definition}[section]
\newtheorem{example}[definition]{Example}
\newtheorem{exercise}{Exercise}[chapter]

\theoremstyle{plain}
\newtheorem{theorem}[definition]{Theorem}
\newtheorem{lemma}[definition]{Lemma}
\newtheorem{corollary}[definition]{Corollary}
\newtheorem{proposition}[definition]{Proposition}

\theoremstyle{remark}
\newtheorem*{remark}{Remark}

\setmathfont{latinmodern-math.otf}
\setmathfont[range=bb]{XITS Math}  

% Common notation
\newcommand{\Z}{\mathbb{Z}}
\newcommand{\Q}{\mathbb{Q}}
\newcommand{\R}{\mathbb{R}}
\newcommand{\C}{\mathbb{C}}
\newcommand{\N}{\mathbb{N}}

% Extra notation
\newcommand{\A}{\mathcal{A}}
\newcommand{\llb}{⟦}
\newcommand{\rrb}{⟧}
\newcommand{\powerset}{\mathbb{P}}

\DeclareMathOperator{\Id}{Id}

\newcommand{\tick}{\ding{51}}
\newcommand{\cross}{\ding{55}}

\title{\textbf{Groups and Monoids in Context}\\[0.5em]
\large An Introduction to Group and Monoid Theory for Undergraduates}
\author{Kazuya Haine}
\date{}

\begin{document}

\maketitle
\tableofcontents

% !TEX root = ../../Groups and Monoids in Context.tex
\setcounter{chapter}{-1}
\chapter{Basics of sets}
% !TEX root = ../../Groups in Context.tex
\section{Sets}
This introductory chapter aims to provide the necessary definitions and intuition about sets required for the rest of the book. The definitions given here will not be those one might find in a formal set theory book, since that is not a very helpful way to introduce ideas. As such, our first definition is as follows:
\begin{definition}[Sets and elements]
    A set is a collection of objects which is only distinguished by what objects it contains. The objects a set contains are called its elements.
\end{definition}
\noindent
This definition begs the question "What is an object?" For the purposes of this book, anything is an object: a carrot, your hat if you're wearing one, the concept of love, etc. Of course, in most cases, we will be using more abstract objects like numbers.

There are a number of ways we can write sets. As is standard from modern algebraic practices, we can simply represent a set by a symbol like $A$, $S$, $x$, etc.
Most of the time, I will use capital letters for sets and lowercase letters for the objects they contain, but there are contexts in which sets contain other sets and for those, this convention somewhat breaks down.

More concretely, we can describe a set by listing its elements. Following mathematical convention, we do this inside curly braces like $\{\text{\emoji{carrot}},\text{\emoji{top-hat}},5, \pi\}$.
\begin{remark}
    Since sets are only distinguished by what objects they contain, there is no such thing as order or repeats in a set, so $\{1, 1\}$ is the same as $\{ 1 \}$, and $\{1,2\}$ is the same as $\{2,1\}$.
    They are however distinguished by whether they contain different things, so $\{1,2\} \ne \{1\}$ because the first contains 2 and the second does not.
\end{remark}
\begin{definition}[Containment]
    If a set $S$ contains an object $x$, we say $x \in S$ read ``$x$ is in $S$'' or ``$x$ is an element of $S$'' or ``$x$ is a member of $S$''. Otherwise, $S$ does not contain $x$, written $x\notin S$.
\end{definition}
\begin{definition}[Subsets]
    A set $A$ is said to be a subset of $B$ if every $a \in A$ is also in $B$. This is written $A \subseteq B$. A subset $A \subseteq B$ is called proper if $A \ne B$ (it doesn't contain \emph{all} the elements of $B$) written $A \subset B$
\end{definition}
\begin{remark}
    The notation for subsets is quite inconsistent among texts. The notation above is one of the main conventions, while the main alternative is using $A \subset B$ to simply mean $A$ is a subset of $B$ and using $\subsetneq$ when specifying proper subsets. Some people may also prefer to use only $\subseteq$ and $\subsetneq$ to be clear at all uses.

    \medskip\noindent
    I have chosen the convention in the definition because it matches the symbols $<$ and $\le$ which will be used frequently when discussing groups.
\end{remark}

\noindent
The main other way of writing sets is by set comprehension
\begin{definition}[Set comprehension]
    If we have a set $S$, we can build a subset of it by restricting to those satisfying certain properties, written $\{ x\in S \mid P(x) \}$ or $\{ x \in S \colon P(x) \}$ where $P(x)$ is a property of $x$.
\end{definition}
\begin{remark}
    We may sometimes also describe a set with a more advanced comprehension which includes an expression on the left hand side of the $\mid$. If this is the case, the right hand side must specify the sets any unknowns belong to (see Example \ref{rational_numbers} for a simple example). This restriction is mostly in place to avoid ill-formed sets\footnote{Russell's paradox is well known and is based on one such ill-formed set: $\{ x \mid x \notin x \}$. The paradox then arises from asking whether this set is an element of itself. I recommend thinking through this question so you can see for yourself the paradox that arises.}.
\end{remark}

\noindent
We now have some basic definitions in place. It would now help to have some common sets in place for us to work with.
\begin{example}[Empty set]
    The empty set is the set which contains nothing, I will normally write it $\varnothing$, but it can also be written $\{\}$ if that is more consistent with the context.
\end{example}
\begin{remark}[The set of everything]
    Due to the restrictions on set comprehension, the most well-used foundation of set theory (ZFC set theory) actually prevents the construction of the set of everything. There are some foundations which do allow it, but they make alternative restrictions to avoid the same issues.
\end{remark}
\begin{example}[Natural numbers]
    The natural numbers are, informally, the counting numbers. They are the numbers 0, 1, 2, 3, \dots. Together, they are written $\N := \{0,1,2,\dots \}$.
\end{example}
\begin{remark}
    Some texts do not include 0 as a natural number (particularly those around number theory), but we have included it in our definition.
\end{remark}
\begin{example}[Integers]
    The integers are all whole numbers, positive and negative. Together, we write them $\Z := \{\dots,-2,-1,0,1,2,\dots\}$.
\end{example}
\begin{example}[Integer intervals]
    For integers $a,b \in \Z$, the set of integers between them (inclusive) is called an integer interval and written $\llb a,b \rrb := \{a,a+1,\dots,b-1,b\}$.
\end{example}
\begin{remark}
    This notation for integer ranges is not universal, but does appear in some texts, and can be quite useful.
\end{remark}
\begin{example}[Rational numbers] \label{rational_numbers}
    The rational numbers are those which can be expressed as a ratio of two integers. We write them $\Q := \left\{\frac{a}{b} \mid a,b \in \Z ,\; b \ne 0 \right\}$.
\end{example}
\begin{example}[Real numbers]
    The real numbers are somewhat difficult to define formally, but they can be thought of as any number on the number line. They are written $\R$ and contain numbers such as $0,1,\sqrt{2}$, and $-\pi$.
\end{example}
\begin{example}[Intervals]
    For real numbers $a,b \in \R$ (with $a < b$), the sets of real numbers between $a$ and $b$ can be written:
    \begin{align*}
        (a,b) &:= \{ x \in \R \mid a < x < b \} \\
        [a,b] &:= \{ x \in \R \mid a \le x \le b \}
    \end{align*}
    These are called the open and closed intervals from $a$ to $b$ respectively.

    \medskip\noindent
    We can also mix open and closed ends to include one end but exclude the other:
    \begin{align*}
        [a,b) &:= \{ x \in \R \mid a \le x < b \} \\
        (a,b] &:= \{ x \in \R \mid a < x \le b \}
    \end{align*}

\end{example}
\begin{example}[Complex numbers]
    The complex numbers are an extension of the real numbers and consist of all numbers of the form $a + ib$ where $a, b \in \R$ and $i$ is a special number satisfying $i^2 = -1$.
    They are written $\C := \{ a + ib \mid a,b \in \R \}$. These will rarely be used in this book, but can occasionally be used to make interesting examples.
\end{example}
\begin{example}[The Latin alphabet]\label{alphabet}
    The letters in the Latin alphabet form a set which will be written:
    \begin{equation*}
        \A := \{\text{a,b,c,d,e,f,g,h,i,j,k,l,m,n,o,p,q,r,s,t,u,v,w,x,y,z}\}.
    \end{equation*}
    I will write them in the standard text font to distinguish them from other algebraic symbols which frequently use an italic font.
\end{example}
\begin{remark}
    This is not a standard notation for the alphabet, but it will be useful in some early examples, so I have included it.
\end{remark}
\begin{example}[Power sets]
    Given a set $S$, we can construct the set of all subsets of $S$ and denote it $\powerset(S)$.
\end{example}
\begin{remark}
    Informally, $\powerset(S) := \{ X \mid X \subseteq S \}$.
    But formally, this set comprehension is not allowed as there is no restriction on the set $X$ belongs to.

    \medskip\noindent
    However, now that power sets have been defined, we can allow such set comprehensions in future by treating $X \subseteq S$ as shorthand for $X \in \powerset(S)$.
\end{remark}

\noindent
3 final operations on sets are particularly useful, so are defined here.

\begin{definition}[Union and intersection]
    If $X$ and $Y$ are sets, we can construct
    \begin{itemize}
        \item the union of $X$ and $Y$ as the set containing all elements of $X$ and all elements of $Y$. This is written $X \cup Y$.
        \item the intersection of $X$ and $Y$ as the set containing all elements in both $X$ and $Y$. This is written $X \cap Y$.
    \end{itemize}
\end{definition}
\begin{definition}[Subtraction]
    If $X$ and $Y$ are sets, then $X \setminus Y := \{ x \in X \mid x \notin Y \}$.
\end{definition}
\begin{remark}
    This is sometimes alternatively written $X - Y$, but in algebraic contexts, this often confuses things, so I will not use this notation.
\end{remark}
\begin{example}
    Let $X = \{1, 3, 5, 7, 9\}$, $Y = \{2, 3, 5, 7\}$. Then:
    \begin{align*}
        X \cup Y &= \{1, 2, 3, 5, 7, 9\}, \\
        X \cap Y &= \{3, 5, 7\}, \\
        X \setminus Y &= \{1, 9\}, \\
        Y \setminus X &= \{2\}.
    \end{align*}
\end{example}
% !TEX root = ../../Groups and Monoids in Context.tex
\section{Tuples and maps}
Now we have a notion of set, we should talk about what sorts of objects will frequently appear besides numbers.
\begin{definition}[products and $n$-tuples]
    If we have a set $A$ and another set $B$, we can form the (cartesian) product of $A$ with $B$:
    \begin{equation*}
        A \times B := \{(a,b) \mid a \in A, b \in B \}.
    \end{equation*}
    Here, $(a,b)$ is an ordered pair, so $(a,b) = (c,d)$ if and only if $a = c$ and $b = d$.

    \medskip\noindent
    Similarly, if we have $n\in \N$ and $n$ sets $A_1,\dots,A_n$, we can form their cartesian product:
    \begin{equation*}
        \prod_{i=1}^n A_i := \{ (a_1,\dots,a_n) \mid \text{$a_i \in A_i$ for each $i \in \llb 1,n \rrb$}\}.
    \end{equation*}
    where each $(a_1,\dots,a_n)$ is an $n$-tuple which is again ordered, so $(a_1,\dots,a_n) = (b_1,\dots,b_n)$ if and only if $a_i = b_i$ for every $i \in \llb 1,n \rrb$.

    \medskip\noindent
    In the case where all $A = A_i$ for every $i \in \llb 1, n \rrb$, we may also denote this set as $A^n$.
\end{definition}
\begin{remark}
    Ordered pairs are just another name for 2-tuples. 3-tuples are often called triples. It is acceptable (but rare) to continue using this naming scheme.
\end{remark}
\begin{remark}
    Notice that for $a,b \in \R$ with $a < b$, the expression $(a,b)$ is ambiguous: it could refer to an open interval, and it could refer to an ordered pair.
    These are unfortunately, both fairly standard notations, and thus, this ambiguity appears a lot.
    It does not generally cause much confusion though, since context almost always makes it clear which way it should be interpreted\footnote{In programming, objects often have a specific type. The same concept can be used to formalise mathematics. This approach largely solves the problem of ambiguities like this, since $(a,b)$ will be specified as either of type Interval or type $\R^2$.}.
\end{remark}

\noindent
Tuples are normally used to define functions, but the definition which uses them is not a helpful definition, it is simply one made out of necessity when people were formalising mathematics with set theory. As such, I present an alternative definition:
\begin{definition}[Maps]
    A map (or function) $f$ is an object with a domain set $A$ and a codomain $B$ which, to each element of $A$ assigns a unique element of $B$. This is denoted $f \colon A \to B$, read ``$f$ is a map from $A$ to $B$''.

    \medskip\noindent
    The unique element of $B$ assigned to a specific $a \in A$ is written $f(a)$, read ``$f$ of $a$'', so $f(a) \in B$ for any $a \in A$.

    \medskip\noindent
    We can also write $f\colon a \mapsto b$ for $a\in A, b\in B$ to say that $f(a) = b$.

    Two functions $f,g \colon A \to B$ are considered equal if $f(a) = g(a)$ for every $a \in A$.
\end{definition}
\begin{remark}
    The unique assignment a function performs does not have to be described by an algorithm to produce it in general.
\end{remark}
\begin{example}
    Following on from the remark, the following is a valid function $\llb 1, 5 \rrb \to \{\text{\emoji{grapes}},\text{\emoji{cat}},\text{\emoji{house}}\}$:
    \begin{align*}
                 1 & \longmapsto \text{\emoji{grapes}} \\
                 2 & \longmapsto \text{\emoji{house}} \\
        f \colon 3 & \longmapsto \text{\emoji{house}} \\
                 4 & \longmapsto \text{\emoji{grapes}} \\
                 5 & \longmapsto \text{\emoji{house}}
    \end{align*}
\end{example}
\begin{example}[Identity]
    For any set $S$, there is a function $\Id_S \colon S \to S$ which does nothing:
    \begin{equation*}
        \Id_S \colon s \mapsto s
    \end{equation*}
    so, $\Id_S(s) = s$.
\end{example}

\begin{definition}[Composition of functions]
    A pair of functions $f \colon A \to B$, $g \colon B \to C$ can be composed to obtain a function $g \circ f \colon A \to C$ given by:
    \begin{equation*}
        g \circ f \colon a \overset{f}{\longmapsto} f(a) \overset{g}{\longmapsto} g(f(a))
    \end{equation*}
\end{definition}
\begin{remark}
    It's important to notice that $g \circ f$ does not mean ``apply $g$, then apply $f$'' but instead the reverse.
    This can confuse people, but this convention is so that the order of functions matches how you would write application without the composition operation:
    \begin{equation*}
        (g \circ f)(a) = g(f(a))
    \end{equation*}
    This is largely due to the convention of how function application is written\footnote{Alternatives conventions include: $f x$, $(x)f$, $x f$. ie. omitting the brackets when unnecessary, reversing the order, and the combination of both. Arguably the last of these is the most natural to use in the contexts this book will cover, but I have chosen the most common notation for reader familiarity.}.
\end{remark}

\begin{definition}[Injections, surjection, and bijections]
    A map $f \colon A \to B$ is called injective if, whenever $f(a_1) = f(a_2)$, $a_1 = a_2$. That is, no two elements of $A$ are mapped to the same element of $B$. Such a map is called an injection.

    \medskip\noindent
    A map $f \colon A \to B$ is called surjective if, for every $b \in B$, there is some $a \in A$ with $f(a) = b$. That is, $f$ maps onto every element of $B$. Such a map is called a surjection.

    \medskip\noindent
    A map $f \colon A \to B$ is called bijective if it is both injective and surjective.
\end{definition}

\begin{definition}
    A map $f \colon A \to B$ is called invertible if there is a map $g \colon B \to A$ such that $g \circ f = \Id_A$ and $f \circ g = \Id_B$. The map $g$ is called the inverse of $f$.
\end{definition}
\begin{remark}
    Formally, this definition does not prove that a function only has 1 inverse if it has one, but Exercise \ref{unique_inverse} shows that inverses are unique if they exist.
\end{remark}

\begin{example}[Identity maps are invertible and bijective]
    To show $\Id_S \colon S \to S$ is invertible, we must find an inverse function $f \colon S \to S$ which reverses the operation of $\Id_S$. But we said that the identity maps do nothing to their elements, so surely if we do nothing to them twice we'll still end up doing nothing? Indeed,
    \begin{equation*}
        \begin{tikzcd}[row sep=0.1em]
            \mathllap{\Id_S \circ \Id_S \;\colon\:} s \arrow[r, maps to, "{\Id_S}"]
                & s \arrow[r, maps to, "{\Id_S}"]
                & s \arrow[d, equal] \\
            \mathllap{\phantom{\Id_S \circ{}}\Id_S \;\colon\:} s \arrow[rr, maps to, "{\Id_S}"']
                & {}
                & s
        \end{tikzcd}
    \end{equation*}
    The two functions agree on their assigned values, and therefore are equal. This covers both left and right compositions of the inverse, so $\Id_S$ is invertible with itself as its inverse.

    \medskip\noindent
    $\Id_S$ is injective if $\Id_S(s_1) = \Id_S(s_2)$ implies that $s_1 = s_2$. This is trivial:
    \begin{equation*}
        s_1 = \Id_S(s_1) = \Id_S(s_2) = s_2
    \end{equation*}
\end{example}
\begin{example}[Addition]
    Addition is normally called a binary operator, which is just a function on pairs:
    \begin{align*}
        +\, \colon \R \times \R &\longrightarrow\:\:\: \R \\
                        (a,b)\; &\longmapsto     a + b
    \end{align*}
    $+$ is not an injection since $1 + 2 = 0 + 3$ for example. ie. $+$ applied to the pairs $(1,2)$ and $(0,3)$ give the same result.

    \medskip\noindent
    $+$ is a surjection since any $r \in \R$ can be obtained as $0 + r$.
    Once we have proven Theorem \ref{bijective_invertible}, we'll see that this means that addition is not invertible.
\end{example}
\begin{example}[Addition by a constant]
    Addition itself is not an injection, but how about addition by a constant? Consider $(c+\phantom{})$ as a function $\R \to \R$ defined by
    \begin{equation*}
        (c+\phantom{}) \colon r \longmapsto c + r
    \end{equation*}
    Now, $(c + \phantom{})$ is injective if whenever $c + r_1 = c + r_2$, we have $r_1 = r_2$. But that's obvious because we can just subtract $c$ from both sides.

    This argument tells us exactly what we should expect the inverse of $(c + \phantom{})$ to be: $(-c + \phantom{})$. Indeed,
    \begin{align*}
        \big((c + \phantom{}) \circ (-c + \phantom{})\big)(r) &= (c + \phantom{})((-c + \phantom{})(r)) \\
        &= (c + \phantom{})(-c + r) \\
        &= c + -c + r \\
        &= r \\
        &= \Id_{\R}(r)
    \end{align*}
    This shows that $(c + \phantom{}) \circ (-c + \phantom{}) = \Id_{\R}$, we also need to show that $(-c + \phantom{}) \circ (c + \phantom{}) = \Id_{\R}$ which follows naturally by a similar chained equality. Thus, $(-c + \phantom{})$ is the inverse map to $(c + \phantom{})$.
    This is the way in which subtraction is the inverse of addition despite the fact that addition itself is not invertible.
\end{example}

\begin{theorem}[bijective if and only if invertible]\label{bijective_invertible}
    Let $f \colon A \to B$ be a map. $f$ is bijective if and only if it is invertible.
\end{theorem}
\begin{proof}
    Suppose $f$ is invertible, then it has inverse $g \colon B \to A$.

    \medskip\noindent
    Injective:

    \noindent
    Let $a_1, a_2 \in A$ be such that $f(a_1) = f(a_2)$, then $g(f(a_1)) = g(f(a_2))$, but $g$ is the inverse of $f$, so
    \begin{equation*}
        a_1 = (g\circ f)(a_1) = g(f(a_1)) = g(f(a_2)) = (g\circ f)(a_2) = a_2.
    \end{equation*}
    Surjective:

    \noindent
    Let $b \in B$, then $g(b) \in A$.
    \begin{equation*}
        f(g(b)) = (f \circ g)(b) = b
    \end{equation*}
    Thus, we have found an element of $A$ which maps onto $b$. This argument did not depend on $b \in B$, so any $b \in B$ has such an element of $A$.

    \medskip\noindent
    Conversely, suppose $f$ is bijective. We wish to construct a function $g \colon B \to A$ which is inverse to $f$.
    Let $b \in B$. Since $f$ is surjective, there is an $a \in A$ such that $f(a) = b$. Since $f$ is injective, only 1 such $a$ exists. This is true for any $b\in B$, so we can define:
    \begin{equation*}
        g \colon b \longmapsto \text{the unique element $a\in A$ such that $f(a) = b$}.
    \end{equation*}
    $g \circ f = \Id_A$:

    \noindent
    $(g \circ f)(a) = g(f(a))$ which is by definition, the unique $a_1 \in A$ such that $f(a_1) = f(a)$. That unique $a_1$ is $a$, so $g(f(a)) = a$.

    \medskip\noindent
    $f \circ g = \Id_B$:

    \noindent
    $(f \circ g)(b) = f(g(b))$. $g(b)$ is an element $a\in A$ with $f(a) = b$, so $f(g(b)) = b$.

    \medskip\noindent
    Together, we see that $g$ is indeed the inverse of $f$, so we are done.
\end{proof}
\begin{remark}
    When we look at the second implication's two pieces, proving $g \circ f = \Id_A$ used the uniqueness of $a_1$ (the injectivity property), while proving $f \circ g = \Id_B$ simply used that the value mapped onto $b$ (the surjectivity property).
    This gives some good indication that injectivity and surjectivity interact quite nicely on specific sides of compositions.
\end{remark}

% !TEX root = ../../Groups in Context.tex
\section{Equivalence and quotients}
The final thing from set theory that will be particularly helpful to know about is the notion of equivalence.
Up until now, I have been quite free about what I mean by equality. 
I mentioned how sets are distinguished only by the elements they contain, which I treated as defining equality for them.
For functions, I was more specific and said that 2 functions are equal if they agree on all outputs.

I didn't dwell on this, but some people might think that two functions defined by different algorithms, but which happen to produce the same values on every input are not truly the same function.
Indeed, they could be considered different objects if they were treated as algorithms or even functions in a programming sense rather than mathematical functions.
This is where the precise language is important: I said that 2 functions are equal when that happens, but equality doesn't have to mean the same in every sense!
This is the essence of what equivalence tries to capture.

To introduce equivalence properly, we need to first introduce relations in general.
\begin{definition}[Relations]
    A (binary) relation is a property defined on a pair of sets $A, B$ which is either satisfied, or not satisfied for any given pair $(a,b) \in A \times B$.
    If the relation is represented by the symbol $\sim$, then $(a, b)$ satisfies $\sim$ is written $a \sim b$.
\end{definition}
\begin{remark}
    The normal way one defines relations is as a generalisation of the function definition that this text has avoided.
    
    One can also think of a relation simply as a binary operator (like $+$, $-$, etc.) which has codomain $\{\text{True}, \text{False}\}$. 
    In this context, $a \sim b$ would be interpreted $\sim(a,b) = \text{True}$. This is non-standard, and so we have avoided this. It is however a good mental model for them.
\end{remark}
\begin{remark}
    There isn't always a way to write ``$a$ is not related to $b$'', however, it is fairly common practice to write this by simply crossing through the relation's symbol:
    \begin{align*}
        \sim \quad &\longleftrightarrow \quad \not\sim \\
        = \quad &\longleftrightarrow \quad \ne \\
        \mid \quad &\longleftrightarrow \quad \nmid
    \end{align*}
\end{remark}
\begin{remark}
    The above definition of relations allows them to relate objects from two different sets. While this is useful in many contexts, we will only be working with relations where the two sets are the same.
    That is, relations will be satisfied by pairs $(x,y) \in A^2$ for a set $A$. As such, from now on a relation on $A$ is a relation on the pair of sets $A$ and $A$.
\end{remark}
\begin{definition}[Symmetry]
    A relation $\sim$ on $A$ is called symmetric if whenever $a \sim b$ (for $a,b\in A$), we also have $b \sim a$.
\end{definition}
\begin{definition}[Reflexivity]
    A relation $\sim$ on $A$ is called reflexive if $a \sim a$ for all $a\in A$.
\end{definition}
\begin{definition}[Transitivity]
    A relation $\sim$ on $A$ is called transitive if whenever $a \sim b$ and $b \sim c$ (for $a,b,c \in A$), we also have $a \sim c$.
\end{definition}
\noindent
These definitions each serve different purposes, and to best familiarise yourself with them, you need to see various examples of relations that satisfy and don't satisfy them. Putting them all together gives us our notion of equivalence.
\begin{definition}[Equivalence]
    A relation $\sim$ on $A$ is called an equivalence relation if it symmetric, reflexive, and transitive. 
\end{definition}

\begin{example}[Equality]
    These are all properties of $=$ on any set, which makes $=$ our main example of an equivalence relation.
    \begin{itemize}
        \item $a = b$ is basically the same thing as $b = a$ when treating it informally.
        \item $a = a$ is something most people would consider a tautology. Indeed, in some areas of logic, it actually satisfies a definition of tautology!
        \item If we know that $a = b$ and $b = c$, then we obviously have $a = c$. This is actually the property that motivates the use of chained equalities like $a = b = c$.
    \end{itemize}
\end{example}
\begin{example}[Ordering]
    There are 4 relations that frequently appear for ordering that aren't equality. Table \ref{order_relation_table} summarises their properties on sets they're defined on.
    \begin{table}[ht]
        \centering
        \begin{tabular}{c|c|c|c}
        Relation & Symmetric & Reflexive & Transitive \\ \hline
        $\le$ & \cross & \tick & \tick \\
        $<$ & \cross & \cross & \tick \\
        $\ge$ & \cross & \tick & \tick \\
        $>$ & \cross & \cross & \tick
        \end{tabular}
        \caption{Ordering relation properties}
        \label{order_relation_table}
    \end{table}
    The fact that $\le$ and $\ge$ have the same properties and $<$ and $>$ do as well is not a coincidence: $a \le b$ is equivalent to $b \ge a$, and similarly for $<$ and $>$.
\end{example}
\begin{definition}[Dual relations]
    If $\triangleright$ and $\triangleleft$ are relations on $A$, they are called dual if $a \triangleleft b$ exactly when $b \triangleright a$.
\end{definition}
\begin{proposition}
    If $\triangleright$ and $\triangleleft$ are dual relations, then,
    \begin{enumerate}[label=(\alph*)]
        \item $\triangleright$ is reflexive if and only if $\triangleleft$ is reflexive,
        \item $\triangleright$ is transitive if and only if $\triangleleft$ is transitive,
        \item $\triangleright$ is symmetric if and only if $\triangleleft$ is symmetric,
    \end{enumerate} 
    Indeed, if they are symmetric, the same pairs satisfy them.
\end{proposition}
\begin{proof}
    Exercise \ref{dual_relations_exercise}.
\end{proof}
\begin{remark}
    While this book will never use category theory, much of what is covered is generalised in category theory. Dual relations are an example of that: 
    they are a special case of dual or opposite categories. Indeed, a part of my motivation for writing this book is to provide a categorical view of groups without requiring readers be familiar with category theory itself.
\end{remark}
\begin{remark}
    We won't go further with dual relations since the final point of the above proposition tells us that equivalence relations are self-dual (this is exactly symmetry).
\end{remark}
\begin{example}[Trivial relation]
    Let $S$ be a set. Let $\boxtimes$ be the relation on $S$ such that no pairs $(x,y) \in S^2$ satisfy it. 
    
    \medskip\noindent
    Symmetric:
    
    \noindent
    Suppose there were a pair $(x,y)\in S^2$ with $x \boxtimes y$ and not $y \boxtimes x$, that would be impossible since no pair satisfies $x \boxtimes y$.

    \medskip\noindent
    Transitive:

    \noindent
    Suppose there were a triple $(x,y,z)\in S^3$ with $x \boxtimes y$ and $y \boxtimes z$, but not $x \boxtimes z$, again, this would be impossible since no pair of $x$ and $y$ exist satisfying $x \boxtimes y$.

    \medskip\noindent
    Possibly reflexive:
    If $S$ is empty, then $\boxtimes$ is vacuously reflexive (like the above two arguments). But if $S$ is not empty, then there is some $x \in S$ and $(x,x)$ does not satisfy $\boxtimes$, so $\boxtimes$ is not reflexive.
\end{example}

\begin{example}[Circle]
    Let $\bigcirc$ be a relation on $\R$ such that $(x,y)\in \R^2$ satisfy $\bigcirc$ precisely when $x^2 + y^2 = 1$. (ie, $(x,y)$ satisfies the relation "the pair lies on the unit circle in the plane centred at the origin".)
    
    \medskip\noindent
    Symmetric:

    \noindent
    Let $(x,y) \in \R^2$. Suppose $x \bigcirc y$, then $x^2 + y^2 = 1$. This means that $y^2 + x^2 = 1$, so $y \bigcirc x$.

    \medskip\noindent
    Not transitive:

    \noindent
    We have $1 \bigcirc 0$ and $0 \bigcirc 1$, but we do not have $1 \bigcirc 1$ since $1^2 + 1^2 = 2 \ne 1$.

    \medskip\noindent
    Not reflexive:

    \noindent
    We do not have $1 \bigcirc 1$.
\end{example}

\begin{example}
    Let $\ddagger$ be a relation on the set $\{\text{\emoji{clubs}},\text{\emoji{diamonds}},\text{\emoji{hearts}},\text{\emoji{spades}}\}$ defined with the relation table in Table \ref{suits_table}:
    \begin{table}[ht!]
        \centering
        \begin{tabular}{c|c|c|c|c}
        $\ddagger$ & \emoji{clubs} & \emoji{diamonds} & \emoji{hearts} & \emoji{spades} \\ \hline
        \emoji{clubs} & \tick & \cross & \tick & \cross \\ \hline
        \emoji{diamonds} & \tick & \tick & \cross & \cross \\ \hline
        \emoji{hearts} & \tick & \cross & \tick & \tick \\ \hline
        \emoji{spades} & \cross & \cross & \tick & \cross 
        \end{tabular}
        \caption{Description of an example relation on suits}
        \label{suits_table}
    \end{table}

    \noindent
    Not symmetric:

    \noindent
    $\text{\emoji{diamonds}} \ddagger \text{\emoji{clubs}}$ but not $\text{\emoji{clubs}} \ddagger \text{\emoji{diamonds}}$.

    \medskip\noindent
    Not transitive:
    
    \noindent
    $\text{\emoji{diamonds}} \ddagger \text{\emoji{clubs}}$, and $\text{\emoji{clubs}} \ddagger \text{\emoji{hearts}}$, but not $\text{\emoji{diamonds}} \ddagger \text{\emoji{hearts}}$.

    \medskip\noindent
    Not reflexive:

    \noindent
    Not $\text{\emoji{spades}} \ddagger \text{\emoji{spades}}$.
\end{example}

\begin{example}[Modular arithmetic]
    Let $n$ be an integer, then let $\equiv_n$ be the relation on $\Z$ where $(a,b) \in \Z^2$ satisfies $\equiv_n$ precisely when $a - b$ is divisible by $n$.

    \medskip\noindent
    Symmetric:
    
    \noindent
    If $a \equiv_n b$, then $a - b = kn$ for some $k \in \Z$. So, $b - a = (-k)n$, so $b - a$ is also divisible by $n$, and $b \equiv_n a$.

    \medskip\noindent
    Transitive:

    \noindent
    If $a \equiv_n b$ and $b \equiv_n c$ for some $a,b,c \in \Z$, then $a - b = k_1n$ for some $k_1 \in \Z$ and $b - c = k_2n$ for some $k_2 \in \Z$. Together, we have $a - c = (a - b) + (b - c) = k_1n + k_2n = (k_1 + k_2)n$, so $a-c$ is divisible by $n$, and $a \equiv_n c$.

    \medskip\noindent
    Reflexive:

    \noindent
    If $a \in \Z$, then $a - a = 0 = 0\times n$, so $a - a$ is divisible by $n$, and $a \equiv_n a$.
\end{example}
\begin{remark}
    This is our first example of an equivalence relation on a set which is not just equality. Indeed, for $n = 3$, we have $1 \equiv_3 7$ despite $1 \ne 7$. As it turns out, for $n = 0$, this does actually just reduce to equality, because multiples of 0 are all identically 0, but all other cases produce non-equality equivalences.
\end{remark}

\begin{definition}[Partitions]
    Let $S$ be a set. A partition of $S$ is a subset $\mathcal{P} \subseteq \powerset(S)$ such that:
    \begin{itemize}
        \item (Covering) for every element $s \in S$, there is some $P \in \mathcal{P}$ with $s \in P$,
        \item (Pairwise disjoint) if $P_1, P_2 \in \mathcal{P}$, then either $P_1 = P_2$ or $P_1 \cap P_2 = \varnothing$.
        \item (Non-empty pieces) every $P\in \mathcal{P}$ is non-empty.
    \end{itemize}
\end{definition}
\begin{remark}
    Mathematical texts have a tendency to label generic objects by the first letter of what they are. This often causes overlaps in notation as seen in the above definition where there are 3 different kinds of the letter P involved. 
    This doesn't happen particularly frequently, but when it does, it can be a bit awkward (particularly when working by hand). The symbols used are completely irrelevant, so in your own work, it is completely acceptable to use different letters.
\end{remark}
\begin{remark}
    Partitions aren't always expressed as a set of subsets, sometimes, they are presented in some other form of collection or just listed.
\end{remark}
\begin{example}
    The sets $\{ x \in \Z \mid \text{$x$ is odd} \}$ and $\{ x \in \Z \mid \text{$x$ is even} \}$ form a partition of $\Z$.
\end{example}
\begin{definition}[Fibres of a map]
    Let $f \colon A \to B$ be a map. For every $b \in B$, we say the fibre of $b$ under $f$ is the pre-image of $b$ under $f$. That is the set
    \begin{equation*}
        \{ a \in A \mid f(a) = b \}
    \end{equation*}
\end{definition}
\begin{remark}
    At this point, it would be nice to show that for any map $f \colon A \to B$, the non-empty fibres of $f$ form a partition of $A$. Indeed, it we could very well at this point, but I'm going to wait so that we can use our next theorem to prove it.
\end{remark}
\begin{definition}[Equivalence classes]
    Let $S$ be a set with an equivalence relation $\sim$ on it. Let $x \in S$, then the $\sim$-equivalence class of $x$ is the set of all elements $y \in S$ for which $x \sim y$. We write this $[x]_{\sim}$ or just $[x]$ if the equivalence relation is clear.
\end{definition}

\begin{theorem}[Equivalences partition]
    If $S$ is a set with an equivalence $\sim$ on it, then the set $\{ [x]_{\sim} \mid x \in S \}$ is a partition of $S$. 
    
    \medbreak\noindent
    Further, if $\mathcal{P}$ is a partition of $S$, then we can define a relation $\sim_{\mathcal{P}}$ defined:
    \begin{equation*}    
        \text{$x \sim_{\mathcal{P}} y$ when $x$ and $y$ lie in the same partition.}
    \end{equation*}
    $\sim_{\mathcal{P}}$ is an equivalence relation whose equivalence classes are exactly the subsets of $S$ in $\mathcal{P}$.
    \label{equivalences_partition}
\end{theorem}
\begin{proof}
    For the first part, every $x \in [x]_{\sim}$ since $x \sim x$ by reflexivity which shows both covering and that all pieces are non-empty. We also need to show that for any pair of $x,y\in S$, the corresponding $[x]_{\sim}, [y]_{\sim}$ are either equal or have no intersection.
    
    \smallskip\noindent
    Suppose they have a non-empty intersection, then there is some $z \in [x]_{\sim} \cap [y]_{\sim}$.
    
    \noindent
    By the definition of equivalence classes, that means $x \sim z$ and $y \sim z$. 
    
    \noindent
    But that means that $z \sim y$ by symmetry. 
    
    \noindent
    Transitivity tells us that $x \sim z \sim y$ means that $x \sim y$, and symmetry gives us $y \sim x$.
    
    \noindent
    So $x \in [y]_{\sim}$ and $y \in [x]_{\sim}$.

    \smallskip\noindent
    Now, suppose that $w \in [x]_{\sim}$. Then $x \sim w$. But we know $y \sim x \sim w$, so $y \sim w$ by transitivity. So $w \in [y]_{\sim}$.
    This shows that $[x]_{\sim} \subseteq [y]_{\sim}$, and a symmetric argument shows the reverse inclusion, so $[x]_{\sim} = [y]_{\sim}$.

    \medskip\noindent
    The second part is exercise \ref{partition_equivalence_exercise}.
\end{proof}

\begin{corollary}[Non-empty fibres partition]
    If $f \colon A \to B$ is a map, then the non-empty fibres of $f$ form a partition of $A$.
\end{corollary}
\begin{proof}
    Let's define an equivalence relation $\sim_f$ on $A$ by $x \sim_f y$ when $f(x) = f(y)$.
    It is easy to check that this is symmetric, transitive, and reflexive (do so!), so this is indeed an equivalence relation.
    Looking at the definitions, the equivalence classes of $\sim_f$ are:
    \begin{equation*}
        [x]_{\sim_f} := \{ y \in A \mid x \sim_f y \} = \{ y \in A \mid f(y) = f(x) \}
    \end{equation*}
    But that's just the fibre of $f(x)$ under $f$.
    Theorem \ref{equivalences_partition} tells us these fibres of the images of elements of $A$ under $f$ partition $A$. These are exactly the non-empty fibres: the fibre of $b$ is non-empty if there is some $a \in A$ such that $f(a) = b$.
\end{proof}

\noindent
We have now defined equivalence relations and shown that they are strongly related to partitions. Now we'll show how we can turn these equivalence relations into equalities.

\begin{definition}[Quotients]
    Let $S$ be a set with an equivalence relation $\sim$ on $S$. Define the quotient of $S$ by $\sim$, written $S/\sim$ by
    \begin{equation*}
        S/\sim \;\;:= \{[x] \mid x \in S \}
    \end{equation*}
\end{definition}
\begin{remark}
    This may not feel any different to things we've done before, but the reason it's important is that we now have $[\cdot] \colon S \to S/\sim$ acting as a map that transforms $x \sim y$ into $[x] = [y]$. This map is also special in the sense that the fibre of $[x] \in S/\sim$ under $[\cdot]$ is exactly $[x]$ itself.
\end{remark}
\noindent
There isn't much more important to say about equivalence relations, partitions, fibres and quotients in the context of sets, but they are very important in the context of groups.
\section*{Exercises}
\begin{exercise}\label{unique_inverse}
    Let $f \colon A \to B$ be a map, and $g_1, g_2 \colon B \to A$ be maps. Suppose both $g_i$ (for $i = 1,2$) satisfy the conditions of being inverse to $f$:
    \begin{equation*}
        g_i \circ f = \Id_A, \quad\quad f \circ g_i = \Id_B.
    \end{equation*}
    Show that $g_1 = g_2$.
\end{exercise}
\begin{exercise}\label{dual_relations_exercise}
    Prove proposition \ref{dual_relations_prop}.
\end{exercise}
\begin{exercise}\label{partitions_equivalence_exercise}
    Complete the proof of proposition \ref{equivalences_partition}
\end{exercise}
\begin{exercise}
    Let $X$ be a set with equivalence relation $\sim$ on it. $f \colon X \to Y$ be a function which is constant on equivalence classes (if $x \sim y$, then $f(x) = f(y)$). Show that there is a function $g \colon X/\sim \;\to Y$ such that $g \circ [\cdot] = f$. Show further that if two such functions $g_1, g_2$ exist, then they must be equal.
\end{exercise}
% !TEX root = ../../Groups and Monoids in Context.tex
\chapter{Groups as Monoids}
% !TEX root = ../../Groups and Monoids in Context.tex
\section{Transformation Groups and Monoids}
We start by introducing groups and monoids in the form the most frequently appear: collections of transformations.
\begin{definition}[Transformation]
    A transformation of a set $S$ is a map $S \to S$. 
\end{definition}
\begin{remark}
    This is not a particularly standard definition of transformation, but the concept of a map from a set to itself is an important concept that deserves a name\footnote{There is standard terminology for this in category theory, but I am restricting such terminology to use for maps between groups and monoids themselves (which will come up in section \ref{sec:homomorphisms}) to avoid confusion.}.
\end{remark}
\begin{definition}[Transformation Monoid]
    A transformation monoid is a collection $M$ of transformations which contains the identity transformation and is closed under composition. That is, $\Id_S \in M$, and whenever $f,g \colon S \to S$ satisfy $f,g \in M$, we also have $f \circ g \in M$.
\end{definition}
\begin{definition}[Transformation Group]
    A transformation group $G$ is a transformation monoid such that every transformation $f \in G$ is invertible with $f^{-1} \in G$.
\end{definition}

\begin{example}[Trivial Transformation Group]
    For any set $S$, the set $\{\Id_S\}$ is a transformation monoid. And since it's only element is the identity element, which is self-inverse, it is also a transformation group. These are the trivial transformation monoids or trivial transformation groups depending on whether we care that they are groups or not.
\end{example}

\begin{example}[Caesar cipher] \textcolor{red}{\boxed{\text{POORLY WRITTEN}}}
    In example \ref{alphabet}, we defined the alphabet $\mathcal{A}$. The Caesar cipher uses a fixed key $k$ between 1 and 25 (normally) and transforms each letter by shifting it forwards in the alphabet $k$ times where each shift either moves it forwards 1 position in the alphabet or for z, sends it back to a.
    Thus, if $c_k \colon \mathcal{A} \to \mathcal{A}$ is the operation the Caesar cipher performs on the alphabet for fixed key $k$, then the set $C_{25} := \{c_k \mid k \in \llb 1, 25 \rrb \}$ is a collection of transformations on $\mathcal{A}$.

    \medskip\noindent
    It's well known (and easy to see), that for each $k \in \llb 1, 25 \rrb$, ${c_k}^{-1} = c_{26-k}$, so every element of this set is invertible, and the inverses are also in the set, but it is neither closed under composition, nor is the identity transformation an element of it, so $C_{25}$ is not a transformation group: $c_1 \circ c_{25} = \Id_{\mathcal{A}} \notin C_{25}$.

    \medskip\noindent
    If we define $c_0 := \Id_{\mathcal{A}}$ and $C_{26} := C_{25} \cup \{c_0\}$, then since this follows modular arithmetic modulo 26, we actually do get that $C_{26}$ (which contains the identity transformation by construction) is closed under composition: whenever $l,m,n \in \llb 0, 25 \rrb$ satisfy $l + m \equiv n \mod 26$, $c_l \circ c_m = c_n$.
\end{example}

\begin{example}[Full Transformation Monoids]
    Let $S$ be a set. The set of all transformations on $S$ is called the full transformation monoid on $S$ (denoted $T_S$) and is indeed a transformtion monoid since the identity transformation on $S$ is a transformation on $S$, and composing transformations on $S$ gives a transformation on $S$.
    If $S$ has at least 2 elements, it is not a transformation group. Indeed, suppose $x,y \in S$ where $x \ne y$. Then, define the transformation $T\colon S \to S$ by
    \begin{equation*}
        T(s) := \begin{cases}
            x & \text{if $s \in \{x,y\}$}, \\
            s & \text{otherwise.}
        \end{cases}
    \end{equation*}
    $T$ is an element of the full transformation monoid of $S$, but $T$ is not invertible since it is not injective: $T(x) = T(y)$ and $x \ne y$.
\end{example}

\begin{example}[Symmetric Groups]
    Let $S$ be a set. The set of all invertible transformations on $S$ is called the symmetric group on $S$, often denoted $\operatorname*{Sym}(S)$. In the cases where $S = \llb 1, n \rrb$, it is often written $S_n$ and called the symmetric group on $n$ symbols. As the names imply, these are transformation groups: The identity transformation is invertible, the composition of invertible transformations is also an invertible transformation, and the inverse of a transformation is also an invertible transformation.
\end{example}

\begin{theorem}[Generated monoids and groups]
    Let $S$ be a set and $T$ be a set of transformations on $S$. Then 
    \begin{enumerate}[label=(\alph*)]
        \item there is a minimal transformation monoid containing $T$ as a subset. It is called the transformation monoid generated by $T$ and denoted $\langle T \rangle_{\text{Mon}}$ or $\langle T \rangle$ if it is clear that we are generating a transformation monoid,
        \item If every element of $T$ is invertible, there is a minimal transformation group containing $T$ as a subset. It is called the transformation group generated by $T$ and is denoted $\langle T \rangle_{\text{Grp}}$ or $\langle T \rangle$ if it is clear that we are generating a transformation group.
    \end{enumerate}
\end{theorem}
\begin{proof}\emph{(a)}

    \smallskip\noindent
    Let $M_0 := \{\Id_S\}$, and define $M_{n+1} := \{ t \circ m \mid t \in T, m \in M_n\}$ inductively. We claim that $\langle T \rangle_{\text{Mon}} = \bigcup_{n = 0}^\infty M_n$.
    To prove this, we must show that $\bigcup_{n = 0}^\infty M_n$ is a transformation monoid and that if $M$ is another transformation monoid containing $T$, then $\bigcup_{n = 0}^\infty M_n \subseteq M$.

    \medskip\noindent
    $\Id_S \in M_0 \subseteq \bigcup_{n = 0}^\infty M_n$, so the identity transformation is an element.

    \medskip\noindent
    If $m_1,m_2 \in \bigcup_{n = 0}^\infty M_n$, then there are $n_1, n_2 \in \mathbb{N}$ with $m_1 \in M_{n_1}, m_2 \in M_{n_2}$. $m_1 = t_1\circ \cdots \circ t_{n_1}$ and $m_2 = t_1' \circ \cdots \circ t_{n_2}'$ where each $t_i, t_j' \in T$. But then $m_1 \circ m_2 = t_1 \circ \cdots \circ t_{n_1} \circ t_1' \circ \cdots \circ t_{n_2}'$, which is clearly an element of $M_{n_1 + n_2}$, and hence an element of $\bigcup_{n = 0}^\infty M_n$.

    \bigskip\noindent
    \emph{(b)}
    
    \smallskip
    Exercise \ref{generated_groups_ex}.
\end{proof}

\begin{example}[Polynomials as functions]
    Let $P$ be the set of polynomials with coefficients in $\mathbb{Z}$. The elements of $P$ can be seen as transformations of $\mathbb{Z}$. The composition of two polynomials is still a polynomial, and the identity function is a polynomial, so $P$ is a transformation monoid.
\end{example}

\begin{example}[Transformation Group of Invertible Elements]
    Let $M$ be a transformation monoid. The set $\{t \in M \mid \text{$t$ is invertible}, t^{-1} \in M \}$ is then a transformation group since we have restricted precisely to the elements which allow it to be a transformation group.
\end{example}

\begin{example}[symmetry groups]
    Let $P_n$ be a set of points on a regular $n$-gon in the plane centred at the origin, then let $D_n$ be the set of rotations and reflections of the plane which fix $P_n$ (sends any point in $P_n$ to another point in $P_n$). $D_n$ is a transformation group on the set $\mathbb{R}^2$ since each rotation and reflection is invertible, and their inverses are themselves rotations or reflections which fix $P_n$, and the identity map on $\mathbb{R}^2$ is a $0\degree$ rotation. The final thing left to check is that composing a rotation with a reflection in either order results in another rotation or reflection (since they both fix $P_n$, their composition does), and some basic linear algebra shows that is true.
\end{example} 


\end{document}